\section{Discussion: Mechanism of Action}
\label{sec:discussion}

While the experimental results quantify the impact of the sponge attack, understanding its underlying mechanism is crucial for designing robust defenses. In this section, we analyze the spectral properties of the adversarial perturbations and the correlation between input features and computational depth.

\subsection{Spectral Stealthiness}
A primary constraint for adversarial attacks in the Smart Grid is stealth. Large, perceptible changes to sensor readings (e.g., wind speed jumping from 5 m/s to 50 m/s) would be immediately flagged by outlier detectors.

We performed a Fast Fourier Transform (FFT) analysis of the generated perturbations $\delta = x_{adv} - x_{benign}$. As shown in Figure \ref{fig:frequency_analysis} (Top Right), the PGD attack concentrates its energy in the high-frequency domain.
\begin{itemize}
    \item \textbf{High-Frequency Injection:} The perturbation spectrum shows distinct peaks at frequencies $f > 0.3$ Hz. In the time domain (Top Left), this manifests as rapid, low-amplitude oscillations rather than global shifts.
    \item \textbf{Sensor Noise Mimicry:} This high-frequency pattern statistically resembles thermal sensor noise or quantization error, making it difficult to distinguish from valid instrument variance using standard low-pass filters.
\end{itemize}

\begin{figure}[t]
    \centering
    \includegraphics[width=\columnwidth]{frequency_analysis.png}
    \caption{\textbf{Spectral Analysis of Perturbations.} (Top) The attack introduces high-frequency oscillations that mimic sensor noise. (Bottom) The spectral power distribution shows that the attack injects the most energy into the 'Temperature' and 'Dewpoint' channels to trigger latency.}
    \label{fig:frequency_analysis}
\end{figure}

\subsection{Feature-Latency Correlation}
To determine which specific input features drive the ACT-LSTM's halting unit, we computed the Pearson correlation coefficient ($r$) between the feature values and the resulting change in ponder steps ($\Delta N$).

Figure \ref{fig:correlations} (Appendix) presents the results:
\begin{itemize}
    \item \textbf{Power ($r \approx +0.40$):} There is a strong positive correlation between the historical Power output and latency. The model "thinks harder" when the grid load is high, likely because the penalty for prediction error is higher in high-load regimes.
    \item \textbf{Wind Speed ($r \approx -0.35$):} Interestingly, Wind Speed shows a negative correlation. This implies that the model halts quickly when wind speeds are high (predictable generation) but hesitates when wind speeds are low or variable (intermittency).
\end{itemize}
The attacker exploits this learned bias by injecting perturbations that artificially lower the perceived wind speed while increasing the perceived load, forcing the model into its "uncertainty zone."

\subsection{The "Logic Trap" Hypothesis}
Combining these findings, we formulate the \textit{Logic Trap Hypothesis}: Dynamic neural networks learn a latent "complexity manifold"—a region of the input space where the prediction is difficult (e.g., sudden weather shifts). The Sponge Attack functions by computing the gradient towards this manifold. Unlike standard integrity attacks that push inputs across a linear decision boundary, sponge attacks push inputs towards the center of high-curvature regions in the loss landscape, trapping the recurrent loop in an extended state of refinement.

\subsection{Model-Size Scaling of Energy via Switching Activity}
The bit-flip results (Table~\ref{tab:bitflip_comparison}) reveal a striking proportionality: the absolute number of switching transitions scales directly with the first-layer weight count. ACT-LSTM (4,608 weights) incurs $\sim$65K baseline flips, DeepAR (18,432 weights) incurs $\sim$262K, and Chronos-v2 (2.1M weights) incurs $\sim$28.4M—a near-linear relationship.

This has an important implication: \textbf{larger models are proportionally more vulnerable to switching-activity energy attacks}, independent of their computational graph structure. While latency-based sponge attacks increase energy only indirectly—and only against dynamic architectures (e.g., ACT-LSTM)—the bit-flip attack vector increases dynamic power dissipation across all architectures equally, including static RNNs and foundation models. As the trend in Smart Grid AI moves towards larger foundation models for improved accuracy, this proportional vulnerability grows in lockstep.

Furthermore, the GA consistently achieves a +25--36\% increase across all model sizes, suggesting that the optimization landscape for bit-flip maximization is smooth and architecture-agnostic. This contrasts sharply with latency attacks, where the GA's effectiveness depends critically on the presence of exploitable control flow (e.g., halting units).

\subsection{Economic Impact Analysis}
Beyond operational stability, sponge attacks impose a tangible financial cost on the operator. We estimate this cost using the "Energy-per-Token" metric.

\subsubsection{Cloud Computing Costs}
For a utility processing $10^6$ inference requests per hour on AWS \textit{p3.2xlarge} instances:
\begin{itemize}
    \item \textbf{Benign Cost:} At 48ms/request, the compute cost is approximately \$11.50 per hour.
    \item \textbf{Attack Cost:} Under the GA sponge attack ($\approx$73.5ms/request), the required compute capacity increases by $53\%$. This forces the auto-scaler to provision additional instances, raising the cost to $\approx \$17.60$ per hour.
    \item \textbf{Annualized Loss:} Over a year, this stealthy latency inflation results in an excess expenditure of $\approx \$53,400$ purely in wasted cloud cycles, not accounting for potential regulatory fines due to delayed reporting.
\end{itemize}

\subsubsection{Battery Replacement Costs}
For edge devices (e.g., IoT sensors with 3000mAh batteries), the 78\% increase in energy consumption from the GA energy attack reduces the device lifespan from 5 years to 2.8 years. This nearly doubles the replacement frequency (maintenance truck rolls), which is the dominant cost factor in AMI operations.
